\documentclass[10pt,a4paper]{amsart}
\usepackage[margin=19mm]{geometry}
\usepackage{amsmath,amssymb,mathtools}
\usepackage[scr=rsfs,cal=euler]{mathalfa}
\usepackage{xfrac}
\usepackage[unicode,hidelinks,bookmarks=false]{hyperref}
\newcommand{\ie}{i.e.\ }
\newcommand{\cf}{cf.\ }
\newcommand{\resp}{respectively\ }
\newcommand{\Subsection}{\subsection}
\providecommand{\nobreakdash}{\nobreak\hbox{-}}
\setlength{\emergencystretch}{2em}
\multlinegap0pt
\usepackage{amssymb}
\usepackage[shortlabels]{enumitem}
\newtheorem{proposition}{Proposition}
\newcommand{\CC}{\mathbb{C}}
\newcommand{\Jac}{\mathrm{rad}}
\newcommand{\LL}{\ell\ell}
\newcommand{\RR}{\mathbb{R}}
\title{Near-Trinions: Complete Classification of\\ Unital Three-Dimensional Real Associative Algebras}
\author{Joel A. Shelton}
\date{Source: arXiv:2606.12773v1 (CC BY 4.0)}
\begin{document}
\maketitle

\noindent\emph{Source and licence.} Near-Trinions: Complete Classification of\\ Unital Three-Dimensional Real Associative Algebras. Version arXiv:2606.12773v1 (CC BY 4.0), distributed by arXiv under the Creative Commons Attribution 4.0 International licence, http://creativecommons.org/licenses/by/4.0/, as recorded in the arXiv licence metadata for this exact version and shown on the abstract page https://arxiv.org/abs/2606.12773v1. Full licence notice: \texttt{licenses/arxiv-2606.12773-CC-BY-4.0.txt}.\par\smallskip

\noindent\textit{Editorial scope.} Counted unit: the article's proposition prop:maxideals (source0.tex lines 350--353). The article's complete original proof of it is delivered with it (source0.tex lines 355--419, one \texttt{proof} environment of 65 lines). No mathematics is written, repaired or re-derived: the statement and the proof are the article's own lines, in the article's own notation, and no retained line is substituted or re-flowed.  No further source-local context is needed: every cross-reference of every retained line resolves inside the two delivered spans.  Nothing was omitted from any retained span, so the package carries no omission record.  No retained line contains a citation, so the wrapper carries no bibliography and no bibliography entry.  No retained line contains a figure environment, an image-inclusion command or drawing code, so no image is delivered and the package declares no asset dependency.  Editorial formatting only, in addition: an AMS \texttt{amsart} preamble that restates the article's own declarations and macros used by the retained lines, namely the environments \texttt{proposition} on separate counters, together with the standard packages those lines need.  The retained environments print in this document as Proposition 1 in order of appearance, whereas the article prints the same items with the numbering of its own full document.  The complete native source of the article as distributed in the arXiv e-print for this exact version is bundled unchanged as \texttt{sources/202647/source0.tex}.
\begin{proposition}\label{prop:maxideals}
  The maximal left ideals of the six near-trinions, with $\Jac(A)$
  confirmed in each case as their intersection, are as follows.
\end{proposition}

\begin{proof}
We work through each class in turn.

\smallskip\noindent
\textbf{(NT1)} $A = \RR^3$, basis $\{e_1, e_2, e_3\}$, $e_ie_j = \delta_{ij}e_i$.
Since $\RR^3$ is semisimple and each $e_i$ generates a minimal ideal
$\RR e_i \cong \RR$, the maximal left ideals are the codimension-one
complementary direct sums:
\[
  \mathfrak{m}_1 = \RR e_2 \oplus \RR e_3,\quad
  \mathfrak{m}_2 = \RR e_1 \oplus \RR e_3,\quad
  \mathfrak{m}_3 = \RR e_1 \oplus \RR e_2.
\]
Their intersection is $0 = \Jac(\RR^3)$.

\smallskip\noindent
\textbf{(NT2)} $A = \RR \times \CC$.
The two minimal direct summands are $\RR \times 0$ and $0 \times \CC$.
The maximal left ideals are their complements:
$\mathfrak{m}_1 = 0 \times \CC$ and $\mathfrak{m}_2 = \RR \times 0$.
Their intersection is $0 = \Jac(\RR \times \CC)$.

\smallskip\noindent
\textbf{(NT3)} $A = \RR \times \RR[\varepsilon]/(\varepsilon^2)$.
The first factor contributes the maximal ideal
$\mathfrak{m}_1 = 0 \times \RR[\varepsilon]/(\varepsilon^2)$.
The second factor $\RR[\varepsilon]/(\varepsilon^2)$ is local with
unique maximal ideal $(\bar\varepsilon)$, contributing
$\mathfrak{m}_2 = \RR \times (\bar\varepsilon)$.
Their intersection is
$\mathfrak{m}_1 \cap \mathfrak{m}_2
 = \{(0, a\bar\varepsilon) : a \in \RR\}
 = 0 \times (\bar\varepsilon) = \Jac(A)$,
one-dimensional, consistent with $d = 1$.

\smallskip\noindent
\textbf{(NT4)} $A = \RR[\varepsilon_1, \varepsilon_2]/(\varepsilon_1^2, \varepsilon_2^2, \varepsilon_1 \varepsilon_2)$.
This algebra is local: an element $a + b\varepsilon_1 + c\varepsilon_2$ is invertible
if and only if $a \neq 0$, with inverse
$a^{-1} - a^{-2}b\varepsilon_1 - a^{-2}c\varepsilon_2$ (using $\varepsilon_a^2 = \varepsilon_1 \varepsilon_2 = 0$).
The unique maximal left ideal is
$\mathfrak{m} = \RR \varepsilon_1 \oplus \RR \varepsilon_2 = \Jac(A)$,
two-dimensional, consistent with $d = 2$.

\smallskip\noindent
\textbf{(NT5)} $A = \RR[x]/(x^3)$.
An element $a_0 + a_1\bar x + a_2\bar x^2$ is invertible iff $a_0 \neq 0$.
The unique maximal ideal is
$\mathfrak{m} = \RR\bar x \oplus \RR\bar x^2 = \Jac(A)$,
two-dimensional with $\Jac(A)^2 = \RR\bar x^2 \neq 0$, so $\LL(A) = 3$.

\smallskip\noindent
\textbf{(NT6)} $A = T_2(\RR)$, basis $\{e_{11}, e_{22}, e_{12}\}$.
An element $ae_{11} + be_{22} + ce_{12}$ is invertible iff $a \neq 0$
and $b \neq 0$.
The two maximal left ideals are:
\[
  \mathfrak{m}_1 = \RR e_{22} \oplus \RR e_{12}
  \quad\text{and}\quad
  \mathfrak{m}_2 = \RR e_{11} \oplus \RR e_{12}.
\]
Their intersection is
$\mathfrak{m}_1 \cap \mathfrak{m}_2 = \RR e_{12} = \Jac(A)$,
one-dimensional, consistent with $d = 1$.
\end{proof}

\end{document}
